\uuid{UIQS}
\exo7id{4120}
\titre{exo7 4120}
\auteur{quercia}
\organisation{exo7}
\datecreate{2010-03-11}
\isIndication{false}
\isCorrection{true}
\chapitre{Equation différentielle}
\sousChapitre{Equations différentielles non linéaires}
\cpge{mpsi-mp2i.fonctions-valeurs-reelles-complexes.b, pcsi.techniques-calcul-fonctions.b, ptsi.techniques-calcul-fonctions.b, mpsi-mp2i.analyse-asymptotique.a, pcsi.analyse-asymptotique.a, ptsi.analyse-asymptotique.a, bcpst1.fonctions-reelles-usuelles.b, mp-mpi.equations-differentielles-lineaires.a, mp-mpi.equations-differentielles-lineaires.b, psi.equations-differentielles, pt.equations-differentielles.a, bcpst2.equations-differentielles-autonomes}
\module{Analyse}
\niveau{L2}
\difficulte{}

\contenu{
\texte{
Soit $y$ une solution maximale de l'équation $y'=x-e^y$.
}
\begin{enumerate}
    \item \question{Montrer que $y$ est décroissante puis croissante.}
\reponse{Régionnement.}
    \item \question{Montrer que $y$ est définie jusqu'en $+\infty$ et que sa courbe représentative
    admet une branche parabolique horizontale.}
\reponse{idem.}
    \item \question{Montrer que $\alpha > -\infty$ et que $y \to +\infty$ lorsque $x\to\alpha^-$.}
\reponse{Pour $x<0$, $y'<-e^y  \Rightarrow  -y'e^{-y}>1  \Rightarrow  x > e^{-y} + C > C$.}
\end{enumerate}
}
